%%%%%%%%%%%%%%%%%%%%%%%%%%%%%%%%%%%%%%%%%%%%%%%%
% E.Bigeard - emmanuel.bigeard@upsti.fr
% http://s2i.bigeard.me
% CC BY-NC-SA 2.0 FR - http://creativecommons.org/licenses/by-nc-sa/2.0/fr/
%%%%%%%%%%%%%%%%%%%%%%%%%%%%%%%%%%%%%%%%%%%%%%%%
\documentclass[11pt]{article}

%%%%%%%%%%%%%%%%%%%%%%%%%%%%%%%%%%%%%%%%%%%%%%%%
% Package UPSTI_Document
%%%%%%%%%%%%%%%%%%%%%%%%%%%%%%%%%%%%%%%%%%%%%%%% 
\usepackage{UPSTI_Document}

%---------------------------------%
% Paramètres du package
%---------------------------------%

% Version du document (pour la compilation)
% 1: Document prof
% 2: Document élève
% 3: Document à publier
\newcommand{\UPSTIidVersionDocument}{3}

% Affichage personnalisé de la classe
\newcommand{\UPSTIclasse}{SNT}

% Type de document
% 0: Custom*				7: Fiche Méthode			14: Document Réponses
% 1: Cours (par défaut)		8: Fiche Synthèse    		15: Programme de colle 
% 2: TD     				9: Formulaire
% 3: TP						10: Memo
% 4: Colle					11: Dossier Technique
% 5: DS						12: Dossier Ressource
% 6: DM						13: Concours Blanc
% * Si on met la valeur 0, il faut décommenter la ligne suivante: 		
%\newcommand{\UPSTItypeDocument}{Custom}
\newcommand{\UPSTIidTypeDocument}{3}

% Titre dans l'en-tête
\newcommand{\UPSTItitreEnTete}{LE LANGAGE CSS}      
\newcommand{\UPSTItitreEnTetePages}{Site web marchand}      
\newcommand{\UPSTIsousTitreEnTete}{Site web marchand}      

% Version du document
\newcommand{\UPSTInumeroVersion}{1.0}
%v1.1 - Ajout de la chaîne fonctionnelle

%----------------------------------------------- 
\UPSTIcompileVars		% "Compile" les variables
%%%%%%%%%%%%%%%%%%%%%%%%%%%%%%%%%%%%%%%%%%%%%%%% 


%%%%%%%%%%%%%%%%%%%%%%%%%%%%%%%%%%%%%%%%%%%%%%%% 
% Début du document
%%%%%%%%%%%%%%%%%%%%%%%%%%%%%%%%%%%%%%%%%%%%%%%% 
\begin{document}
\UPSTIbuildPage

\UPSTIobjectif{Sur cette deuxième activité vous allez modifier le code CSS du site, chaque élève sera responsable d'une page contenant plusieurs feuille de style. 
}

\competencetable{Compétences à valider}{
\competence{Etudier et modifier une page HTML simple}
}

\section*{CHOISIR LES COULEURS}

\activiteressources
{Définir une palette de couleurs pour l'ensemble du site}
{
\task Choisir une palette de couleur pour le site ;
\task Relever le code HTML des différentes couleurs à appliquer au site web ;
\task Télécharger le dossier "styles" et le dézipper dans votre espace réseau avec le dossier de l'activité précédente.
}
{documentation}
{\url{https://htmlcolorcodes.com/fr/noms-de-couleur/}}


\section*{AJOUTER LES LIENS CSS}

\activiteressources
<<<<<<< HEAD
{Ajouter dans le code HTML de votre page les différents liens vers les fichiers CSS.}
=======
{Ajouter dans le code HTML de votre page, les différents liens vers les fichiers CSS.}
>>>>>>> Activite1
{
\task Rechercher pour ma page les fichiers CSS correspondant ;
\task Ajouter les liens vers les fichiers CSS dans le code HTML ;
}
{documentation}
<<<<<<< HEAD
{NotePad}
=======
{NotePad ++}
>>>>>>> Activite1

\section*{MODIFIER LE STYLE DU SITE}

\activiteressources
{Mettre en place le style de votre page web.}
{%
\task Modifier les couleurs pour qu'elles correspondent à la palette de couleur choisie ;
\task Adapter si besoin la position des éléments et des boutons.
}
{documentation}
<<<<<<< HEAD
{NotePad}
=======
{NotePad ++}
>>>>>>> Activite1

\section*{TESTER LE STYLE COMPLET}

\activiteressources
{Tester le style du site}
{
\task Sur un PC du groupe, copier les différentes pages modifiées et les feuilles de styles CSS ;
\task Visualiser le style de votre site
}
{documentation}
<<<<<<< HEAD
{NotePad}
=======
{NotePad ++}
>>>>>>> Activite1



\end{document}
